\documentclass[11pt,a4paper,twocolumn]{article}

\usepackage[utf8]{inputenc}
\usepackage[T1]{fontenc}
\usepackage[spanish]{babel}
\usepackage[top=2.5cm,bottom=2.5cm,left=2cm,right=2cm,columnsep=0.6cm]{geometry}
\usepackage{xcolor}
\usepackage{booktabs}
\usepackage{array}
\usepackage{tabularx}
\usepackage{enumitem}
\usepackage{fancyhdr}
\usepackage{lastpage}
\usepackage{amsmath}
\usepackage{amssymb}
\usepackage{parskip}
\usepackage{titlesec}
\usepackage{titling}
\usepackage{abstract}
\usepackage{caption}
\usepackage{listings}
\usepackage[colorlinks=true,linkcolor=teal,urlcolor=teal,citecolor=teal]{hyperref}

% ── Colores ────────────────────────────────────────────────────
\definecolor{accent}{HTML}{0d9488}
\definecolor{textmuted}{HTML}{4a5568}
\definecolor{warncolor}{HTML}{b45309}
\definecolor{critcolor}{HTML}{b91c1c}
\definecolor{codebg}{HTML}{f8f9fa}
\definecolor{codefg}{HTML}{2d3748}

% ── Tipografía de secciones ────────────────────────────────────
\titleformat{\section}
  {\large\bfseries\color{accent}}
  {\thesection.}{0.5em}{}[\vspace{-4pt}\color{accent}\rule{\linewidth}{0.6pt}]

\titleformat{\subsection}
  {\normalsize\bfseries\color{accent!80!black}}
  {\thesubsection.}{0.5em}{}

\titleformat{\subsubsection}
  {\normalsize\itshape\color{textmuted}}
  {\thesubsubsection.}{0.5em}{}

% ── Cabecera / pie ─────────────────────────────────────────────
\setlength{\headheight}{14pt}
\pagestyle{fancy}
\fancyhf{}
\fancyhead[L]{\small\color{textmuted}\textit{Artificial World --- Motor de Civilizaciones Emergentes}}
\fancyhead[R]{\small\color{textmuted}2026-03-08}
\fancyfoot[C]{\small\color{textmuted}P\'agina \thepage\ / \pageref{LastPage}}
\renewcommand{\headrulewidth}{0.4pt}

% ── Listings (código) ──────────────────────────────────────────
\lstset{
  backgroundcolor=\color{codebg},
  basicstyle=\footnotesize\ttfamily\color{codefg},
  breaklines=true,
  frame=single,
  framerule=0.4pt,
  rulecolor=\color{accent!40},
  xleftmargin=6pt,
  xrightmargin=6pt,
  aboveskip=6pt,
  belowskip=6pt,
}

% ── Abstract ──────────────────────────────────────────────────
\renewcommand{\abstractnamefont}{\normalfont\bfseries\color{accent}}
\renewcommand{\abstracttextfont}{\small\itshape}
\setlength{\absleftindent}{0pt}
\setlength{\absrightindent}{0pt}

% ── Tablas ────────────────────────────────────────────────────
\renewcommand{\arraystretch}{1.25}

% ──────────────────────────────────────────────────────────────
\begin{document}

% ── PORTADA ───────────────────────────────────────────────────
\twocolumn[{%
\begin{@twocolumnfalse}
  \vspace*{0.5cm}
  {\color{accent}\rule{\linewidth}{2pt}}\\[0.6em]
  {\LARGE\bfseries\color{accent}
    Artificial World: Un Motor de Simulaci\'on de\\ Civilizaciones
    Emergentes con Agentes Aut\'onomos Basados en Utilidad}\\[0.5em]
  {\color{accent}\rule{\linewidth}{0.6pt}}\\[0.4em]
  {\normalsize
    \textbf{Autor:} Cosigein SL\quad
    \textbf{Fecha:} 2026-03-08\quad
    \textbf{Repositorio:} \texttt{artificial-word}\quad
    \textbf{Versi\'on:} Sesi\'on fundacional, semilla~42, 200~ticks
  }\\[0.4em]
  {\color{accent}\rule{\linewidth}{2pt}}\\[0.8em]

  \begin{abstract}
    Presentamos \textbf{Artificial World}, un motor de simulaci\'on de
    vida artificial 2D orientado a la generaci\'on de civilizaciones
    emergentes con memoria, h\'eroes, refugios y comunidades.
    El sistema implementa agentes aut\'onomos cuya toma de decisiones
    se fundamenta en una funci\'on de utilidad multidimensional que
    integra estado fisiol\'ogico, rasgos de personalidad, memoria
    espacial y social, relaciones inter-agente, directivas externas y
    reglas de autonom\'ia de supervivencia.
    Describimos la arquitectura del motor de decisi\'on, la estructura
    del entorno simulado, el protocolo de sesi\'on fundacional
    reproducible y los resultados observados en una sesi\'on can\'onica
    de 200~ticks con 8~entidades sobre un mapa de $60\times60$ celdas.
    Los resultados muestran un estado de tensi\'on sostenida con
    supervivencia total de la poblaci\'on pero con hambre cr\'itica en
    el 87.5\% de los agentes al final del experimento, evidenciando la
    presi\'on evolutiva del entorno y la insuficiencia del mecanismo de
    b\'usqueda de recursos en condiciones de escasez prolongada.
    El sistema es reproducible mediante semilla fija y est\'a disponible
    como motor Python open source con persistencia SQLite.\\[4pt]
    \textbf{Palabras clave:} vida artificial, agentes aut\'onomos,
    funci\'on de utilidad, civilizaciones emergentes, simulaci\'on 2D,
    comportamiento emergente, memoria social, modo sombra.
  \end{abstract}
  \vspace{0.6em}
\end{@twocolumnfalse}
}]

% ── 1. INTRODUCCIÓN ───────────────────────────────────────────
\section{Introducci\'on}

La simulaci\'on de sistemas de vida artificial ha sido un campo f\'ertil
desde los trabajos fundacionales de Conway (1970) con el Juego de la
Vida~\cite{Conway1970} y los ecosistemas de Reynolds (1987) con
Boids~\cite{Reynolds1987}.
Sin embargo, la mayor\'ia de los sistemas cl\'asicos priorizan la
simplicidad de las reglas locales sobre la riqueza del estado interno
de los agentes.

Artificial World parte de una premisa distinta: \textbf{los agentes
deben tener estado interno complejo, memoria persistente y capacidad de
decidir mediante ponderaci\'on de utilidades}, de forma que el
comportamiento colectivo emergente sea consecuencia de decisiones
individualmente racionales bajo restricci\'on.

El objetivo del sistema no es la optimizaci\'on de una m\'etrica global,
sino la generaci\'on de \textbf{historia emergente}: secuencias de
eventos que den lugar a cr\'onicas, memorias y finalmente a
civilizaciones con identidad propia. La unidad fundamental no es el
agente individual sino el \textbf{refugio} --- un nodo de supervivencia
colectiva alrededor del cual se organiza la comunidad.

Este trabajo documenta: (1)~arquitectura del motor de decisi\'on por
utilidad; (2)~entorno simulado y par\'ametros; (3)~protocolo de sesi\'on
fundacional reproducible; (4)~resultados de la sesi\'on can\'onica
(semilla~42, 200~ticks); (5)~an\'alisis de comportamiento emergente;
(6)~limitaciones y trabajo futuro.

% ── 2. TRABAJO RELACIONADO ────────────────────────────────────
\section{Trabajo Relacionado}

\subsection{Vida Artificial y Agentes Aut\'onomos}

Los sistemas de vida artificial basados en agentes (Langton,
1989~\cite{Langton1989}) se dividen en dos familias:
\textbf{sistemas de reglas locales simples} (aut\'omatas celulares,
Boids) y \textbf{sistemas de agentes con estado interno} (Creatures,
The Sims, Dwarf Fortress~\cite{Adams2006}).
Artificial World pertenece a la segunda familia, con \'enfasis en la
\textbf{arquitectura de decisi\'on por utilidad} y en la
\textbf{persistencia de memoria individual}.

\subsection{Utilidad en Agentes Deliberativos}

La teor\'ia de la utilidad esperada~\cite{VonNeumann1944} ha sido
adaptada extensamente para sistemas multiagente.
La arquitectura se inspira en los sistemas \textbf{BDI}
(Rao \& Georgeff, 1991~\cite{Rao1991}), pero reemplaza la selecci\'on
de intenciones por una funci\'on de puntuaci\'on continua que permite
graduaci\'on natural entre opciones.

\subsection{Sistemas de Refugio y Comunidad}

La noci\'on de refugio como unidad de organizaci\'on social tiene
antecedentes en Dwarf Fortress, RimWorld~\cite{Sylvester2013} y Caves
of Qud. Artificial World formaliza el refugio como \textbf{entidad de
primer orden} con estado propio, en lugar de tratarlo como mero
artefacto del mapa.

% ── 3. ARQUITECTURA ───────────────────────────────────────────
\section{Arquitectura del Sistema}

\subsection{Visi\'on General}

El sistema se compone de las capas mostradas en el
Listado~\ref{lst:arch}.
Dispone de dos motores independientes:

\begin{itemize}[leftmargin=1em, itemsep=2pt, topsep=2pt]
  \item \textbf{Motor Python} (Python~3.11+ / pygame): simulaci\'on
        can\'onica, alta fidelidad; persistencia SQLite
        (\texttt{mundo\_artificial.db}).
  \item \textbf{Motor JS} (Node.js / Express): demo web, HeroRefuge,
        API; persistencia SQLite (\texttt{audit\_simulacion.db}).
\end{itemize}

La simulaci\'on avanza por ticks discretos. En cada tick, cada entidad
ejecuta el ciclo: \textit{Percepci\'on} $\to$ \textit{Generaci\'on de
candidatos} $\to$ \textit{Puntuaci\'on} $\to$ \textit{Selecci\'on}.

\begin{lstlisting}[label=lst:arch, caption={\'Arbol de m\'odulos principal}]
principal.py
|-- nucleo/simulacion.py
|   |-- agentes/motor_decision.py
|   |-- acciones/        (13 tipos)
|   |-- sistemas/        (Watchdog, Modo Sombra)
|   `-- mundo/mapa.py    (grid 60x60)
|-- sistemas/cronica_fundacional.py
|-- backend/src/         (API Node.js)
`-- frontend/src/        (React + Tailwind)
\end{lstlisting}

\subsection{Funci\'on de Utilidad}

La puntuaci\'on de una acci\'on $a$ para el agente $i$ en el tick $t$ es:

\begin{equation}
  U(a,i,t) = U_{\text{base}}(a) + \sum_{k} \delta_k(a,i,t) + \varepsilon
  \label{eq:utility}
\end{equation}

donde $U_{\text{base}}(a)$ es la utilidad intr\'inseca de la acci\'on
(Tabla~\ref{tab:ubase}), $\delta_k$ son modificadores contextuales
(Tabla~\ref{tab:modificadores}), y
$\varepsilon \sim \mathcal{U}(-0.08, 0.08)$ es ruido determinista
por semilla.

\begin{table}[h]
  \centering
  \caption{Utilidades base por acci\'on}
  \label{tab:ubase}
  \small
  \begin{tabular}{lc}
    \toprule
    \textbf{Acci\'on} & $U_{\text{base}}$ \\
    \midrule
    \texttt{comer}             & 0.65 \\
    \texttt{recoger\_comida}   & 0.55 \\
    \texttt{explorar}          & 0.50 \\
    \texttt{huir}              & 0.40 \\
    \texttt{descansar}         & 0.30 \\
    \texttt{mover}             & 0.30 \\
    \texttt{recoger\_material} & 0.30 \\
    \texttt{ir\_refugio}       & 0.25 \\
    \texttt{evitar}            & 0.25 \\
    \texttt{compartir}         & 0.20 \\
    \texttt{robar}             & 0.20 \\
    \texttt{seguir}            & 0.20 \\
    \texttt{construir}         & 0.18 \\
    \bottomrule
  \end{tabular}
\end{table}

\begin{table}[h]
  \centering
  \caption{Modificadores de la funci\'on de utilidad}
  \label{tab:modificadores}
  \small
  \begin{tabular}{p{1.4cm}p{3.6cm}}
    \toprule
    \textbf{Modificador} & \textbf{Efecto} \\
    \midrule
    hambre    & $\uparrow$ comer/recoger cuando hambre $>$ umbral \\
    energ\'ia & $\uparrow$ descansar cuando energ\'ia $<0.6$ \\
    riesgo    & $\uparrow$ huir/evitar cuando riesgo $>0.3$ \\
    rasgo     & Modifica preferencias por arquetipo \\
    oscilaci\'on & $-0.40$ sobre retorno a celda anterior \\
    memoria   & Bonus hacia comida recordada \\
    relaciones & Confianza/hostilidad/miedo con cercanos \\
    directivas & \'Ordenes externas con prioridad alta \\
    autonom\'ia & Override supervivencia sobre directivas \\
    \bottomrule
  \end{tabular}
\end{table}

La capa de autonom\'ia act\'ua como \textbf{override de emergencia}:
si $\text{hambre}>0.85$, se penaliza exploraci\'on; si
$\text{energ\'ia}<0.15$, se penaliza movimiento y robo; si
$\text{riesgo}>0.8$, se potencia la huida sobre cualquier directiva.

\subsection{Cat\'alogo de Acciones}

El sistema implementa 13~tipos de acci\'on, cada uno con condici\'on de
viabilidad. Si ninguna es viable, el sistema garantiza fallback a
\texttt{DESCANSAR}.

\subsection{Rasgos de Agente}

Los agentes sociales tienen 5~arquetipos: \textit{Cooperativo},
\textit{Neutral}, \textit{Agresivo}, \textit{Explorador} y
\textit{Oportunista}.
Los agentes de tipo \texttt{gato} tienen su propia tabla de 5~rasgos:
\textit{Curioso}, \textit{Apegado}, \textit{Independiente},
\textit{Territorial} y \textit{Oportunista}.

\subsection{Memoria de Agente}

Cada agente mantiene dos registros:

\begin{itemize}[leftmargin=1em, itemsep=2pt, topsep=2pt]
  \item \textbf{Memoria espacial}: posiciones de recursos observados
        recientemente; alimenta directamente el modificador
        \textit{memoria} de la ec.~(\ref{eq:utility}).
  \item \textbf{Memoria social}: relaciones con otros agentes en tres
        dimensiones: \texttt{confianza}, \texttt{hostilidad},
        \texttt{miedo}.
\end{itemize}

\subsection{Modo Sombra y Watchdog}

El \textbf{Modo Sombra} es una simulaci\'on especular paralela que
permite evaluar contraf\'acticos sin alterar el mundo observable.
El \textbf{Watchdog} monitoriza condiciones an\'omalas y emite alertas
estructuradas con nivel (\texttt{WARN}, \texttt{CRITICAL}), c\'odigo
sem\'antico, tick de emisi\'on y entidad afectada.

% ── 4. ENTORNO EXPERIMENTAL ───────────────────────────────────
\section{Entorno Experimental}

\subsection{Par\'ametros de la Sesi\'on Can\'onica}

\begin{table}[h]
  \centering
  \caption{Par\'ametros de la sesi\'on fundacional}
  \label{tab:params}
  \small
  \begin{tabular}{ll}
    \toprule
    \textbf{Par\'ametro} & \textbf{Valor} \\
    \midrule
    Semilla aleatoria         & 42 \\
    Dimensiones del mapa      & $60 \times 60$ celdas \\
    Entidades iniciales       & 8 \\
    Ticks ejecutados          & 200 \\
    Fundador                  & Tryndamere \\
    Refugio fundador          & Refugio Fundador \\
    Semilla de civilizaci\'on & default \\
    Timestamp                 & 2026-03-08T12:53:48 \\
    \bottomrule
  \end{tabular}
\end{table}

\subsection{Composici\'on de la Poblaci\'on}

\begin{table}[h]
  \centering
  \caption{Estado final de las 8 entidades (tick~200)}
  \label{tab:poblacion}
  \small
  \begin{tabular}{llcc}
    \toprule
    \textbf{Nombre} & \textbf{Tipo} & \textbf{H} & \textbf{E} \\
    \midrule
    Ana         & social & \textcolor{warncolor}{0.68} & 0.38 \\
    Bruno       & social & \textcolor{critcolor}{1.00} & 0.36 \\
    Clara       & social & \textcolor{critcolor}{1.00} & 0.26 \\
    David       & social & \textcolor{critcolor}{0.90} & \textcolor{warncolor}{0.18} \\
    Eva         & social & \textcolor{critcolor}{1.00} & 0.40 \\
    F\'elix     & social & \textcolor{critcolor}{1.00} & 0.38 \\
    Amiguisimo  & gato   & \textcolor{critcolor}{1.00} & 0.30 \\
    Tryndamere  & social & \textcolor{critcolor}{1.00} & 0.40 \\
    \bottomrule
    \multicolumn{4}{l}{\footnotesize H = hambre [0=saciado, 1=cr\'itico]; E = energ\'ia [0=agotado, 1=pleno]}
  \end{tabular}
\end{table}

La poblaci\'on incluye 7~entidades de tipo \texttt{social} y
1~de tipo \texttt{gato} (Amiguisimo). Todas mantienen
$\text{salud}=1.0$ al final, indicando que ninguna alcanz\'o el
estado de muerte.

% ── 5. RESULTADOS ─────────────────────────────────────────────
\section{Resultados}

\subsection{M\'etricas de Acci\'on Agregadas}

\begin{table}[h]
  \centering
  \caption{Distribución de acciones en 200 ticks}
  \label{tab:acciones}
  \small
  \begin{tabular}{lrc}
    \toprule
    \textbf{Acci\'on} & \textbf{Ejec.} & \textbf{Proporci\'on} \\
    \midrule
    \texttt{SEGUIR}          & 998 & 62.4\% \\
    \texttt{DESCANSAR}       & 492 & 30.8\% \\
    \texttt{RECOGER\_RECURSO} &  55 &  3.4\% \\
    \texttt{COMER}           &  55 &  3.4\% \\
    \midrule
    \textbf{Total}           & \textbf{1.600} & \textbf{100\%} \\
    \bottomrule
  \end{tabular}
\end{table}

La dominancia de \texttt{SEGUIR} (62.4\%) revela un patr\'on emergente:
ante la escasez de recursos, los agentes con rasgo social priorizan la
cohesi\'on grupal sobre la b\'usqueda activa individual. Este resultado
es consistente con la hip\'otesis de que el modificador de relaciones
sociales supera al modificador de hambre en los estadios medios del
experimento.

\subsection{Estado Final de la Poblaci\'on}

\begin{itemize}[leftmargin=1em, itemsep=2pt, topsep=2pt]
  \item \textbf{Veredicto:} \textcolor{warncolor}{\textsc{tensi\'on}}
  \item \textbf{Entidades vivas:} 8/8 (supervivencia total)
  \item \textbf{Hambre m\'axima:} 1.00 (cr\'itica)
  \item \textbf{Hambre promedio:} 0.9475
  \item \textbf{Energ\'ia m\'inima:} 0.18
  \item \textbf{Energ\'ia promedio:} 0.3325
  \item \textbf{Alertas emitidas:} 65
\end{itemize}

\subsection{An\'alisis de Alertas Watchdog}

Las 65~alertas se distribuyen en dos c\'odigos sem\'anticos:
\texttt{HAMBRE\_SIN\_COMIDA\_DISPONIBLE} (WARN) y
\texttt{HAMBRE\_CRITICA\_SIN\_RESPUESTA} (CRITICAL).
Se concentran en los ticks~189 y~199, indicando presi\'on cr\'itica
sostenida en la fase final (ticks 180--200).
Los agentes afectados (Tryndamere, Amiguisimo, Eva, Clara, Bruno,
F\'elix) est\'an ubicados en la mitad izquierda del mapa ($x<15$),
zona con menor densidad de recursos en esta configuraci\'on de semilla.

\subsection{Dispersi\'on Espacial}

Los vectores de posici\'on final revelan fragmentaci\'on en tres
clusters:

\begin{itemize}[leftmargin=1em, itemsep=2pt, topsep=2pt]
  \item \textbf{Cluster NW} ($x<15$, $y>40$): Tryndamere (2,51),
        F\'elix (0,52), Amiguisimo (12,42).
  \item \textbf{Cluster SW} ($x<30$, $y<25$): David (10,22),
        Eva (13,5), Clara (25,5).
  \item \textbf{Posici\'on aislada}: Bruno (0,34), Ana (47,21).
\end{itemize}

% ── 6. DISCUSIÓN ──────────────────────────────────────────────
\section{Discusi\'on}

\subsection{Emergencia del Comportamiento Social}

El resultado m\'as significativo es la dominancia de \texttt{SEGUIR}
(62.4\%) en un contexto de escasez severa. Este comportamiento no fue
codificado como respuesta a la escasez, sino que emerge de la
interacci\'on entre: (1)~el modificador de relaciones (confianza alta
$\to$ bonus \texttt{SEGUIR}); (2)~la penalizaci\'on de autonom\'ia
sobre exploraci\'on cuando $\text{hambre}>0.85$; y (3)~la ausencia de
comida que hace inviables \texttt{COMER} y \texttt{RECOGER} la mayor
parte del tiempo.

Esto produce una \textbf{trampa din\'amica}: los agentes cohesionados
mantienen relaciones de alta confianza que refuerzan el seguimiento
mutuo, pero esto reduce la exploraci\'on necesaria para encontrar
recursos.

\subsection{El Problema del Agotamiento Zonal}

La concentraci\'on de alertas en el cluster oeste sugiere que la
semilla~42 genera un mapa con gradiente de recursos que desfavorece
las zonas de borde oeste. Los agentes all\'i ubicados sufren hambre
cr\'itica cr\'onica por \textbf{agotamiento zonal} --- han consumido
los recursos locales y el modificador anti-oscilaci\'on ($-0.40$)
dificulta el retroceso eficiente.

\subsection{El Caso An\'omalo: Ana}

Ana es la \'unica entidad con hambre sub-cr\'itica ($H=0.68$), ubicada
en $(47,21)$ --- la posici\'on m\'as alejada del grupo. Su aislamiento
espacial, no intencional sino resultado de la exploraci\'on individual,
result\'o adaptativo: sin competencia zonal, tuvo acceso exclusivo a
recursos del cuadrante NE.

Este resultado sugiere que, bajo presi\'on de escasez, la
\textbf{divergencia individual de la cohesi\'on grupal} puede ser una
estrategia emergentemente superior, aunque el sistema no la codifica
como tal.

\subsection{Reproducibilidad}

La sesi\'on es completamente reproducible mediante semilla fija:

\begin{lstlisting}[language=bash, caption={Comando de reproducci\'on}]
python cronica_fundacional.py \
  --seed 42 --ticks 200 \
  --founder "Tryndamere" \
  --refuge "Refugio Fundador"
\end{lstlisting}

% ── 7. ARQUITECTURA CONCEPTUAL ────────────────────────────────
\section{Modelo de Civilizaciones Vivas}

\subsection{Entidades de Dominio}

\begin{table}[h]
  \centering
  \caption{Entidades del modelo de dominio}
  \label{tab:dominio}
  \small
  \begin{tabular}{p{1.8cm}p{3.8cm}}
    \toprule
    \textbf{Entidad} & \textbf{Rol} \\
    \midrule
    \texttt{World}            & Contenedor global \\
    \texttt{CivilizationSeed} & Arquetipo fundacional \\
    \texttt{Refuge}           & Unidad base de supervivencia \\
    \texttt{Hero}             & Agente hist\'orico \\
    \texttt{Community}        & Agrupaci\'on viva \\
    \texttt{MemoryEntry}      & Registro de memoria \\
    \texttt{HistoricalRecord} & Evento hist\'orico \\
    \texttt{Territory}        & Lectura estrat\'egica 2D \\
    \bottomrule
  \end{tabular}
\end{table}

\subsection{Flujo Fundador}

El flujo m\'inimo viable es:
\textit{Semilla} $\to$ \textit{Refugio} $\to$ \textit{Civilizaci\'on
naciente}. El usuario elige una \texttt{CivilizationSeed}, nombra al
h\'eroe fundador y al refugio; el sistema crea entidades, comunidad y
cr\'onica fundacional.

\subsection{Estrategia 2D / 3D}

\begin{itemize}[leftmargin=1em, itemsep=2pt, topsep=2pt]
  \item \textbf{2D} (implementada): verdad sist\'emica --- mapa, grid,
        rutas, nodos, recursos, refugios, influencia, fronteras.
  \item \textbf{3D} (roadmap): encarnaci\'on futura --- h\'eroes de gran
        presencia, refugios encarnados, eventos hist\'oricos visuales.
\end{itemize}

% ── 8. VERIFICACIÓN ───────────────────────────────────────────
\section{Verificaci\'on y Tests}

El sistema cuenta con 11~suites de tests Python y 45~tests en
backend Node/Vitest, con un total verificado de \textbf{68+~tests}
pasando en CI en cada push.
Las suites cubren: ciclo tick, acciones, decisi\'on, cr\'onica
fundacional, Modo Sombra, combate, interacciones sociales, Watchdog
e integraci\'on producci\'on completa.

Para ejecutar:

\begin{lstlisting}[language=bash, caption={Runner de tests}]
$env:SDL_VIDEODRIVER="dummy"
$env:SDL_AUDIODRIVER="dummy"
python pruebas/run_tests_produccion.py
\end{lstlisting}

% ── 9. LIMITACIONES Y TRABAJO FUTURO ──────────────────────────
\section{Limitaciones y Trabajo Futuro}

\subsection{Limitaciones Actuales}

\begin{enumerate}[leftmargin=1.2em, itemsep=2pt, topsep=2pt]
  \item \textbf{Ausencia de integraci\'on Python/JS}: los dos motores
        son independientes sin comunicaci\'on en tiempo real.
  \item \textbf{Escasez de hitos narrativos}: solo 2~hitos en 200~ticks
        (inicio y cierre); eventos intermedios no se registraron en esta
        sesi\'on.
  \item \textbf{Comunidad no simulada}: \texttt{Community} existe como
        contrato de datos, pero no como simulaci\'on social rica.
  \item \textbf{Diplomacia y linajes}: civilizaciones completas con
        guerras y migraciones permanecen en roadmap.
  \item \textbf{3D no implementada}: la capa de encarnaci\'on visual no
        existe a\'un como runtime.
  \item \textbf{Capturas visuales}: la cr\'onica a\'un no incluye
        screenshots autom\'aticos de ticks clave.
\end{enumerate}

\subsection{Trabajo Futuro}

Las pr\'oximas l\'ineas de desarrollo son: enriquecimiento de
\texttt{Refuge} y \texttt{Community} con normas sociales emergentes;
historia profunda con m\'as eventos alimentando \texttt{MemoryEntry};
mecanismo de migraci\'on grupal ante agotamiento zonal; implementaci\'on
completa de las 7~semillas de civilizaci\'on (tribu, tecn\'ocrata,
espiritual, guerrero, comerciante, paranoica, decadente); e
integraci\'on Python/JS con puente en tiempo real.

% ── 10. CONCLUSIONES ──────────────────────────────────────────
\section{Conclusiones}

Artificial World demuestra que un motor de simulaci\'on relativamente
compacto (arquitectura Python 2D, 13~tipos de acci\'on, funci\'on de
utilidad con 9~modificadores) es capaz de producir comportamientos
emergentes no triviales y verificables. La sesi\'on can\'onica
(semilla~42, 200~ticks) revela:

\begin{itemize}[leftmargin=1em, itemsep=2pt, topsep=2pt]
  \item \textbf{Supervivencia total} bajo presi\'on de escasez severa.
  \item \textbf{Emergencia de cohesi\'on social} como respuesta a la
        escasez, con el efecto paradójico de reducir la exploraci\'on
        individual.
  \item \textbf{Divergencia individual exitosa} como estrategia no
        codificada (caso Ana).
  \item \textbf{Agotamiento zonal} como causa primaria de hambre
        cr\'itica, no fallo algorítmico.
  \item \textbf{Reproducibilidad completa} mediante semilla fija y
        protocolo documentado.
\end{itemize}

El sistema establece las bases para un motor creador de mundos compacto
y reutilizable. La tesis de producto ---
\textbf{civilizaciones vivas con refugios, h\'eroes y memoria} ---
es t\'ecnicamente defendible en su estado actual en la dimensi\'on 2D,
con la capa de encarnaci\'on 3D expl\'icitamente marcada como roadmap.

\smallskip
\begin{center}
  \textit{``Empieza con un refugio. Elige una semilla.\\
  Mira nacer tu civilizaci\'on.''}
\end{center}

% ── REFERENCIAS ───────────────────────────────────────────────
\begin{thebibliography}{9}
  \bibitem{Adams2006}
    Adams, T.\ (2006). \textit{Dwarf Fortress}. Bay~12 Games.

  \bibitem{Conway1970}
    Conway, J.\ H.\ (1970). The Game of Life.
    \textit{Scientific American}, 223(4), 4--10.

  \bibitem{Langton1989}
    Langton, C.\ G.\ (1989). Artificial Life.
    \textit{Artificial Life}, SFI Studies in the Sciences of Complexity,
    1--47.

  \bibitem{Rao1991}
    Rao, A.\ S., \& Georgeff, M.\ P.\ (1991).
    Modeling Rational Agents within a BDI-Architecture.
    \textit{KR '91}, 473--484.

  \bibitem{Reynolds1987}
    Reynolds, C.\ W.\ (1987). Flocks, Herds, and Schools:
    A Distributed Behavioral Model.
    \textit{SIGGRAPH Computer Graphics}, 21(4), 25--34.

  \bibitem{Sylvester2013}
    Sylvester, T.\ (2013). \textit{RimWorld}. Ludeon Studios.

  \bibitem{VonNeumann1944}
    Von Neumann, J., \& Morgenstern, O.\ (1944).
    \textit{Theory of Games and Economic Behavior}.
    Princeton University Press.
\end{thebibliography}

\end{document}
